\section{第 2 章 相关技术与理论基础}

\subsection{2.1 Flask 与前后端集成式 Web 架构}

本系统使用的是 Flask 作为后端的核心框架。相比于更重型的全栈框架，Flask的优势在于结构轻量、路由组织清晰、与模板页面和独立接口结合灵活，适合本项目这样“页面展示+数据接口+分析逻辑并存”的综合型应用。系统中的主页、多模式预报页、探空分析页、历史气候页和农业气象页均由 Flask 路由统一组织，而数据接口如 /api/advanced-forecast/multi-model、/upperair/data、/apply-ml-correction 等则承担前后端交互和异步数据更新功能。

该架构的业务意义就在于，它既保持了服务端模板渲染适合页面级整合展示的优点，又使关键模块可以独立地通过异步接口来更新数据，从而避免整页刷新造成的等待以及耦合问题。多模式预报、探空分析等都必须在页面加载之后再请求数据，这样的模板首屏+接口补充模式更适合气象信息平台的实时交互。

从系统组织角度来说，Flask 对于将路由层、服务层、分析层、模型层进行拆分是十分友好的。路由层主要是负责页面入口和接口调度，服务层主要是对外部数据进行获取和缓存组织，分析层主要是对历史气候、探空参数和农业指标进行计算，模型层主要是加载并调用机器学习误差校正模型。分层并不追求复杂的架构形式，但是可以很好地满足目前系统扩展的需求。

\subsection{2.2 气象数据来源与特点}

本系统使用的数据来源具有明显的多源异构特征，不同模块依赖的数据在时间尺度、空间尺度和更新机制上均存在差异，因此理解这些数据源的特点，是后续系统实现和结果解释的基础。

在数值预报方面，系统主要依赖 Open-Meteo 提供的开放接口获取 ECMWF、GFS、ICON 等模式结果。其优点在于调用便捷、格式统一、适合快速构建多模式对比能力。对本系统而言，Open-Meteo 不只是一个“拿天气数据”的接口，更是多模式预报分析和 ECMWF 后处理应用的基础输入源。

历史气象和农业积温用 ERA5-Land 再分析资料作为长期历史背景数据源。再分析资料的时间连续性好、历史资料丰富、字段统一，适合用作历史趋势分析、长期气候背景描述和积温这类长期积分型指标的估计。尽管再分析数据不是地面单站观测，但是在长期累积和区域背景分析方面有较好的稳定性以及可用性。

系统用怀俄明大学在线探空资料服务做高空观测数据的来源。该数据源可以给出标准探空时次的层结表以及指数资料，是做Skew-T分析、大气稳定度判断和对流潜势诊断的基础。与地面观测数据相比，探空资料可以更直接地反映出垂直方向上的温湿风结构，在系统中起到专业的诊断作用。

观测真值和机器学习误差校正方面，目前系统首先统一到杭州国家站/馒头山口径，如果更高的权威和稳定的本地小时观测获取有困难，那么阶段性使用Meteostat 58457作为当前最好的可用的统一真值源。虽然存在一定的工程折中，但是比训练、验证、运行时混用不同站点口径更合理，也更容易保证模型评价结果的可解释性。

\subsection{2.3 数据处理与分析支撑技术}

本系统在数据处理上主要使用了Pandas和NumPy这两个库。Pandas主要用在时间序列数据解析、缺测值清洗、按日或者按小时聚合、统计汇总、结构化输出上，NumPy更多地用来做数值计算、百分位统计、趋势计算、特征构造等。对气象系统来说，组合起来很适合处理时间字段多、缺测情况复杂、统计需求频繁的数据形态。

历史气候模块的年统计、月统计、同月同日样本提取、EW指数计算、季节事件筛选等都是用Pandas和NumPy完成的；农业模块中积温累计、适宜度评分、净有效降水等指标的计算也都基于这些基础处理能力之上。因此可以认为这两个工具是系统分析层最根本的数据处理基础。

另外系统在运行的时候还会加入简化的缓存机制来减少高频的数据请求造成的性能损失。多模式预报结果、模块状态以及部分提示信息都会被缓存起来，在短时间内不会对外部数据源进行重复请求。该种方式虽然不是复杂的分布式缓存方案，但是对目前规模的毕业设计系统来说，在保证时效性的同时也改善了响应体验。

\subsection{2.4 可视化与页面交互技术}

气象平台的结果表达高度依赖图形可视化，因此本系统在不同模块中采用了不同的图表技术组合。

在多模式预报和历史气候模块中，系统主要使用 Plotly.js 以及后端 Plotly / Plotly Express 生成交互式时间序列图、月度对比图、趋势图和分布图。Plotly 的优势在于交互能力较强，适合展示温度、降水、风速等连续变化过程，也方便在同一坐标系内对比多个模式或多个年份。

在探空分析模块中，系统主要使用 Matplotlib 与 MetPy。其中MetPy给出气象分析常用的绘图、热力学计算的支持，适合T-lnP图、Skew-T图、风矢量图等专业的图形绘制，Matplotlib是底层的绘图支持，保证探空图输出具有灵活性以及专业性。对于该模块来说，选择Matplotlib+MetPy比通用Web图表库更适合，因为探空图本身就是一种专业性的气象图形。

农业模块页面布局主要是用 Bootstrap，专题页环境曲线用 Chart.js 绘制。bootstrap适合于组织多卡片、多区域的信息页面，Chart.js适合绘制较轻量的专题曲线图。前端交互部分用JavaScript实现异步数据请求和局部刷新，农业页面中AI建议刷新、探空页面的图形切换等都属于此类。

该种“不同的模块使用不同的图形技术”是符合系统实际需要的。也就是说，本系统并没有因为统一技术栈而强行使用一种图表库，而是根据业务场景选择更加适合的表达方式。

\subsection{2.5 常用气象分析方法基础}

本系统虽然不以理论创新为目标，但仍然用到了若干具有明确气象业务背景的分析方法。

在多模式预报分析中，核心思想是通过同一时空位置下不同模式结果的并列展示，帮助用户判断天气过程强度、转折时间和趋势演变的一致性或分歧程度。这种方法本质上属于数值预报结果的比较分析，强调“同一事件、不同模式”的横向判断。

在探空分析中，系统使用 CAPE、CIN、K 指数、Lifted Index、整层可降水量和垂直风切变等指标来辅助判断大气不稳定条件和潜在风险。其中，Skew-T 或 T-lnP 图是分析温湿层结、抬升凝结和对流发展环境的经典图形工具，也是专业气象分析中的重要基础。

历史气候模块里，系统除了保留传统的原始极值排序之外，又加入了以历史同月同日的经验百分位为基础的EW指数，用来衡量某一日天气事件在长期历史背景下所具有的相对异常程度。对低温使用反向百分位、连续事件过程去重、四季分组展示，都是为了使极端事件识别更符合气候分析的业务逻辑。

农业模块中积温计算、适宜度评分属于比较典型的分析手段。积温用基点温度以上日平均有效热量的累计来表示作物发育进程，适宜度评分综合考虑温度、水分、极端天气的影响，对当前环境是否有利于作物生长做出定量的描述。虽然这些方法使用了简化的实现方式，但是已经可以满足农业管理提示的需要。

\subsection{2.6 机器学习误差校正基础}

机器学习误差校正在本系统中的定位，不是构建一个独立研究型算法平台，而是作为 ECMWF 预报链路的后处理增强层。当前系统采用 scikit-learn 中的随机森林模型完成温度、湿度、风速和降水的偏差订正，其中温度、湿度、风速使用回归模型，降水采用“先分类、后回归”的两阶段方案。

这一设计与业务问题本身是匹配的。温度、湿度和风速属于连续变量，更适合直接回归订正；降水则既包含“下不下”的分类问题，也包含“下多少”的量级问题，因此拆分建模更符合实际情况。特征方面，系统主要使用原始预报值、小时、月份、星期、昼夜属性、时效长度 lead\_hours、时效分桶以及周期性时间编码等变量，目的不是追求复杂模型，而是构建与运行时接口兼容、同时具备基本解释力的订正方案。

系统用issue\_time组织的滚动验证和最终留出测试来代替随机拆分所有目标时刻的方法。这样可以减少信息泄露，更接近真实业务中用过去起报样本预测未来起报样本的使用场景。本论文中这一部分的重要意义不单是模型本身，也是它体现出来的系统结果可信度以及工程闭环的提高。

\subsection{2.7 机器学习误差校正相关算法}

为了更清楚地说明本系统中误差校正模块的建模思路，本文对V1和V2两版方案中涉及到的主要算法进行简要说明。前者主要是用随机森林做回归和分类，后者也使用梯度提升树中的CatBoost做对比实验。

（1）随机森林(Random Forest)。随机森林属于一种基于决策树的集成学习方法，它主要依靠Bootstrap重采样来创建多个训练子集，然后单独训练多棵决策树，最后在预测时将各个树的结果加以整合。回归任务中模型一般输出各个树的结果平均值，分类任务用多数投票法得到最终类别。该方法抗过拟合能力强，能处理非线性关系，对异常值和局部噪声不敏感，适合当前样本量小、业务特征强的气象误差校正场景，所以本文V1模型使用了该算法。

（2）梯度提升树（CatBoost）。梯度提升树是通过不断训练新的树来拟合前一轮模型的残差，从而不断减小整体预测误差的一种方法。CatBoost是梯度提升树的高效工程实现，训练速度快、类别特征支持好、泛化能力强。在本文中，V2实验模型试图用CatBoost对ECMWF误差校正任务进行重训，以验证在统一站点口径、统一issue\_time规则、严格时间隔离评估条件下，boosting类模型能否得到比V1更好的结果。

